\documentclass[11pt,a4paper]{article}

\usepackage[utf8]{inputenc}
\usepackage[T1]{fontenc}
\usepackage[spanish,es-tabla]{babel}
\usepackage[margin=2.35cm]{geometry}
\usepackage{amsmath}
\usepackage{amssymb}
\usepackage{bm}
\usepackage{graphicx}
\usepackage{booktabs}
\usepackage{array}
\usepackage{float}
\usepackage{xcolor}
\usepackage{microtype}
\microtypesetup{expansion=false}
\usepackage{hyperref}

\hypersetup{
    colorlinks=true,
    linkcolor=blue!60!black,
    citecolor=blue!60!black,
    urlcolor=blue!60!black
}

\graphicspath{{artifacts/figures/}{artifacts/figures/generated/}}

\title{\textbf{Sistema de apoyo al análisis inicial de casos ante el TEDH:}\\
\large recuperación de precedentes y predicción explicable de artículos jurídicos con análisis de deriva temporal}

\author{
    Jordi Blasco Lozano\\
    \small Grado en Ingeniería en Inteligencia Artificial\\
    \small Escuela Politécnica Superior, Universidad de Alicante\\
    \small Aprendizaje Avanzado --- Curso 2025/2026
}

\date{}

\begin{document}
\maketitle

\begin{abstract}
Este trabajo estudia ECtHR Task B de LexGLUE como un problema de apoyo al análisis inicial de casos ante el Tribunal Europeo de Derechos Humanos. Dado un caso descrito por sus hechos, el sistema propuesto sugiere artículos candidatos del Convenio Europeo de Derechos Humanos, recupera precedentes similares y proporciona señales de auditoría sobre sus predicciones. El objetivo no es automatizar decisiones judiciales ni predecir sentencias, sino priorizar hipótesis de trabajo revisables por un analista jurídico. El dataset contiene 11000 casos, divididos en 9000 de entrenamiento, 1000 de validación y 1000 de test, con 10 artículos, 15991 pares positivos caso-artículo y 262580 párrafos almacenados. La metodología combina una base SQLite trazable, representación TF-IDF con unigramas y bigramas, clasificación One-vs-Rest, ajuste de umbrales por etiqueta, recuperación de precedentes mediante similitud coseno y análisis de explicabilidad global/local. Los resultados en test alcanzan macro-F1 de 0.672, micro-F1 de 0.737 y Hamming loss de 0.076. En retrieval, la variante TF-IDF obtiene Recall@10 de 0.968 sobre una muestra reproducible de 250 consultas. Finalmente, se analiza deriva de etiquetas y patrones de error, destacando las implicaciones prácticas de falsos negativos, falsos positivos, anclaje por precedentes y límites de las explicaciones basadas en términos.
\end{abstract}

\noindent\textbf{Palabras clave:} NLP jurídico, clasificación multietiqueta, LexGLUE, ECtHR, explicabilidad, drift temporal, recuperación de precedentes.

\tableofcontents
\newpage

\section{Introducción}

El análisis de jurisprudencia del Tribunal Europeo de Derechos Humanos (TEDH) exige leer documentos extensos, identificar artículos potencialmente implicados y localizar precedentes relevantes. En una fase inicial de estudio de un caso, un abogado, investigador o analista jurídico no necesita una ``decisión automática'', sino una priorización razonable de líneas de investigación: qué artículos revisar, qué casos previos pueden ser comparables y qué señales textuales han llevado al sistema a sugerir una hipótesis.

Este trabajo aborda ese escenario mediante aprendizaje automático aplicado a ECtHR Task B de LexGLUE. La tarea se formula como clasificación multietiqueta de artículos del Convenio Europeo de Derechos Humanos a partir de los hechos del caso. A diferencia del encuadre de predicción de sentencias, el sistema se presenta como apoyo al \textit{issue spotting} jurídico y a la recuperación de precedentes. Esta distinción es central: una predicción de artículo no es razonamiento jurídico, no resuelve el caso y no debe utilizarse sin supervisión humana.

El proyecto se centra en dos problemáticas propias de sistemas de aprendizaje automático en dominios reales. La primera es la explicabilidad: en derecho, una sugerencia opaca tiene valor limitado y puede inducir confianza indebida. La segunda es la deriva temporal: la distribución de casos, artículos y vocabulario puede cambiar con el tiempo, y esos cambios pueden reflejar evolución jurisprudencial, cambios en litigación o sesgos de muestreo. La privacidad se discute como requisito de despliegue responsable, pero no constituye el eje experimental porque el benchmark utilizado es público.

La contribución del trabajo es un pipeline reproducible que combina ingesta en SQLite, EDA, modelado multietiqueta con TF-IDF y One-vs-Rest, ajuste de umbrales, recuperación de precedentes, explicaciones globales/locales y análisis de drift y errores. La Figura~\ref{fig:overview} resume el sistema completo.

\begin{figure}[H]
\centering
\includegraphics[width=0.94\textwidth]{fig01_system_overview.pdf}
\caption{Overview del sistema. El flujo transforma un caso jurídico bruto en artículos candidatos, precedentes recuperados y señales de auditoría. El sistema apoya al analista jurídico y no decide sentencias.}
\label{fig:overview}
\end{figure}

\section{Trabajo relacionado}
\label{sec:related}

\subsection{Benchmarks y modelos de NLP jurídico}

LexGLUE propone un conjunto de tareas para evaluar comprensión de lenguaje jurídico en diferentes dominios, incluyendo casos del TEDH~\cite{chalkidis2022lexglue}. Este tipo de benchmark permite comparar modelos, pero también obliga a interpretar con cuidado qué significa cada tarea. ECtHR Task B no requiere decidir el resultado final de un proceso, sino identificar artículos asociados a los hechos. Por ello, en este trabajo se evita el marco de ``predicción automática de sentencias'' y se usa una formulación de apoyo al análisis inicial.

La literatura sobre modelos jurídicos contextualizados, como LEGAL-BERT~\cite{chalkidis2020legalbert}, muestra que adaptar representaciones al dominio legal puede mejorar tareas de clasificación. Sin embargo, el objetivo aquí no es maximizar rendimiento con arquitecturas grandes, sino construir un sistema reproducible, auditable y fácil de discutir. TF-IDF y modelos lineales siguen siendo útiles como línea base fuerte cuando se necesita inspeccionar pesos, umbrales y errores.

\subsection{Predicción jurídica y alcance del problema}

El trabajo de Aletras et al.~\cite{aletras2016predicting} popularizó la predicción de decisiones del TEDH mediante análisis textual. Esa línea ha sido influyente, pero también discutida por el riesgo de confundir correlaciones textuales con razonamiento legal. En este proyecto se adopta una posición más conservadora: el modelo no predice sentencias ni sustituye interpretación jurídica, sino que sugiere artículos candidatos y casos similares para acelerar una revisión humana.

La recuperación de precedentes se relaciona con sistemas de búsqueda jurídica y \textit{case-based reasoning}. En esta entrega se usa una definición operacional de relevancia: un precedente recuperado es relevante si comparte al menos un artículo verdadero con la consulta. Esta aproximación es reproducible y barata, pero no equivale a una evaluación experta de similitud doctrinal.

\subsection{Explicabilidad y deriva}

LIME~\cite{ribeiro2016lime} explica predicciones locales perturbando entradas y ajustando un modelo interpretable alrededor del caso. En modelos lineales también pueden inspeccionarse coeficientes globales. Ambas estrategias son útiles para auditoría, pero no deben interpretarse como causalidad jurídica: un término puede correlacionar con un artículo sin representar por sí mismo una razón normativa.

La deriva de datos y de concepto es una preocupación clásica en aprendizaje automático real. Gama et al.~\cite{gama2014survey} revisan adaptación ante \textit{concept drift}, mientras que Moreno-Torres et al.~\cite{moreno2012datasetshift} sistematizan tipos de cambio de distribución. En derecho, el drift puede ser especialmente ambiguo: puede indicar degradación estadística, pero también evolución jurisprudencial o cambios en el tipo de litigios. Como inspiración visual para representar fallos y estados del sistema se tuvo en cuenta el estilo de análisis de fallos de \textit{The Illusion of Thinking}~\cite{shojaee2025illusion}, sin usarlo como fundamento jurídico.

\section{Descripción del problema y los datos}
\label{sec:datos}

\subsection{Definición formal}

Dado un conjunto de casos \(D=\{(x_i,y_i)\}_{i=1}^{n}\), cada entrada \(x_i\) es el texto completo de los hechos de un caso del TEDH. La salida \(y_i\) es un vector binario:
\[
y_i \in \{0,1\}^{K}, \qquad K=10,
\]
donde cada dimensión representa un artículo del Convenio. El objetivo es aprender una función:
\[
f: \mathcal{X} \rightarrow [0,1]^K
\]
que asigne un score por artículo. La decisión final se obtiene con umbrales por etiqueta:
\[
\hat{y}_{ij} = \mathbb{1}(f_j(x_i) \geq \tau_j).
\]
El ajuste de \(\tau_j\) se realiza en validación para mejorar el equilibrio por etiqueta. La partición de test se reserva para evaluación final.

\subsection{Dataset}

El dataset usado es LexGLUE / ECtHR Task B. Tras la ingesta completa, se materializa en SQLite con 11000 casos y 262580 párrafos. La Tabla~\ref{tab:dataset} resume sus características.

\begin{table}[H]
\centering
\caption{Características principales del dataset ECtHR Task B utilizado.}
\label{tab:dataset}
\small
\begin{tabular}{@{}lr@{}}
\toprule
\textbf{Característica} & \textbf{Valor}\\
\midrule
Casos de entrenamiento & 9000\\
Casos de validación & 1000\\
Casos de test & 1000\\
Casos totales & 11000\\
Artículos / etiquetas & 10\\
Pares positivos caso-artículo & 15991\\
Media de artículos por caso etiquetado & 1.477\\
Máximo de artículos en un caso & 7\\
Párrafos almacenados & 262580\\
\bottomrule
\end{tabular}
\end{table}

Las etiquetas son los artículos \(2,3,5,6,8,9,10,11,14\) y \(P1\)-\(1\). La tarea es multietiqueta: un mismo caso puede activar varios artículos. Además, hay fuerte desbalanceo; algunos artículos aparecen en miles de casos y otros son minoritarios.

\subsection{Análisis exploratorio}

La Figura~\ref{fig:eda} sintetiza el EDA. El split de entrenamiento domina en tamaño, los documentos tienen longitudes muy variables y la frecuencia de artículos muestra una cola larga. La matriz de coocurrencia confirma que ciertos artículos aparecen juntos con frecuencia, por lo que el sistema debe permitir múltiples etiquetas simultáneas. Esta evidencia justifica usar macro-F1 junto a micro-F1: micro-F1 resume rendimiento agregado, mientras que macro-F1 penaliza peor el bajo rendimiento en etiquetas raras.

\begin{figure}[H]
\centering
\includegraphics[width=\textwidth]{fig03_eda_multipanel.pdf}
\caption{EDA multipanel. El dataset combina texto largo, desbalanceo, coocurrencias entre artículos y naturaleza multietiqueta.}
\label{fig:eda}
\end{figure}

\subsection{Estructura de datos}

El pipeline almacena el dataset en \path{data/interim/metadata.db}. Las tablas principales son \texttt{cases}, \texttt{case\_paragraphs}, \texttt{articles}, \texttt{case\_labels}, \texttt{predictions}, \texttt{experiment\_runs} y \texttt{explanations}. Esta estructura separa texto, párrafos, etiquetas, predicciones y metadatos de ejecución. La separación facilita auditoría: una predicción puede vincularse al caso, al modelo, a sus scores y a la configuración experimental.

\begin{figure}[H]
\centering
\includegraphics[width=0.96\textwidth]{fig02_state_transition_pipeline.pdf}
\caption{Estados del flujo. Esquema compacto generado desde notebook: la entrada textual se transforma en representación estructurada, scores de modelo y salida jurídica revisable.}
\label{fig:states}
\end{figure}

\section{Metodología}
\label{sec:metodologia}

\subsection{Representación y clasificación}

La representación principal es TF-IDF con unigramas y bigramas, \texttt{min\_df=2}, \texttt{max\_df=0.95}, \texttt{sublinear\_tf=True} y un máximo de 60000 características. El vectorizador se ajusta únicamente sobre \texttt{train} para evitar fuga de información. Con la matriz \(X\) y la matriz binaria \(Y\), se entrenan modelos One-vs-Rest: Logistic Regression como modelo principal y Linear SVM como comparación suplementaria.

El modelo principal usa scores probabilísticos de regresión logística y umbrales específicos por etiqueta ajustados en validación. Esta decisión responde al desbalanceo: un umbral global de 0.5 puede ser demasiado conservador para artículos poco frecuentes o demasiado permisivo para etiquetas dominantes.

\begin{figure}[H]
\centering
\includegraphics[width=0.96\textwidth]{fig04_modeling_and_thresholds.pdf}
\caption{Modelado multietiqueta. Se muestran el pipeline TF-IDF + OVR, las métricas principales por split, la distribución de scores por artículo y los umbrales ajustados.}
\label{fig:modeling}
\end{figure}

\subsection{Retrieval de precedentes}

El módulo de recuperación usa los casos de \texttt{train} como índice de precedentes. Se evalúan dos variantes: TF-IDF con similitud coseno y TF-IDF reducido con TruncatedSVD seguido de coseno. Para mantener el notebook interactivo y reproducible, el índice se construye sobre una muestra de 2500 casos de train y se evalúa sobre 250 consultas por split (\texttt{validation} y \texttt{test}).

La relevancia se define de forma operacional: un caso recuperado es relevante si comparte al menos un artículo verdadero con la consulta. Las métricas usadas son Recall@5, Recall@10 y nDCG@10. Esta evaluación no sustituye una revisión jurídica experta del precedente; solo mide si el ranking recupera casos con solapamiento de artículos.

\subsection{Explicabilidad}

La explicabilidad se aborda en dos niveles. Globalmente, se extraen los términos con mayores coeficientes positivos y negativos por artículo en el modelo lineal. Localmente, se generan explicaciones LIME para cinco casos de test. Además, se incluye una prueba mínima de fidelidad: retirar un término clave y observar el cambio de score. La prueba no valida causalidad, pero ayuda a comprobar si una señal destacada influye de forma medible en el modelo.

\subsection{Drift y errores}

El drift se mide comparando la distribución de artículos de train con validation y test. Se usan divergencia Jensen-Shannon, distancia L1 y similitud coseno. También se calculan deltas de macro-F1, micro-F1 y Hamming loss respecto a validation. El análisis de errores calcula verdaderos positivos, falsos positivos y falsos negativos por caso, y agrupa patrones frecuentes FP/FN por split.

\section{Experimentos y resultados}
\label{sec:resultados}

\subsection{Clasificación multietiqueta}

La Tabla~\ref{tab:classification} contiene los resultados principales. El modelo alcanza en test macro-F1 de 0.672, micro-F1 de 0.737 y Hamming loss de 0.076. La diferencia entre train y validation/test muestra que el problema no es trivial y que existe sobreajuste moderado o cambio de distribución entre particiones.

\begin{table}[H]
\centering
\caption{Resultados principales de clasificación. Fuente: tabla final en artifacts/metrics.}
\label{tab:classification}
\small
\begin{tabular}{@{}llrrrr@{}}
\toprule
\textbf{Modelo} & \textbf{Split} & \textbf{Casos} & \textbf{Macro-F1} & \textbf{Micro-F1} & \textbf{Hamming}\\
\midrule
TF-IDF LogReg OVR + umbrales & train & 9000 & 0.816 & 0.876 & 0.037\\
TF-IDF LogReg OVR + umbrales & validation & 1000 & 0.654 & 0.758 & 0.069\\
TF-IDF LogReg OVR + umbrales & test & 1000 & 0.672 & 0.737 & 0.076\\
\bottomrule
\end{tabular}
\end{table}

El notebook genera además una comparación suplementaria con baseline, Logistic Regression a umbral 0.5, Linear SVM a umbral 0.5 y Logistic Regression con umbrales ajustados. Para el cuerpo principal se reporta únicamente el modelo final seleccionado, evitando mezclar resultados exploratorios con la tabla principal.

\subsection{Retrieval}

La Figura~\ref{fig:retrieval} y la Tabla~\ref{tab:retrieval} resumen el módulo de precedentes. En test, TF-IDF coseno logra Recall@10 de 0.968 y nDCG@10 de 0.868; TF-IDF+SVD obtiene nDCG@10 de 0.870 y Recall@10 de 0.956. La diferencia es pequeña y no debe sobreinterpretarse porque la evaluación es muestral.

\begin{figure}[H]
\centering
\includegraphics[width=0.96\textwidth]{fig05_retrieval_precedents.pdf}
\caption{Retrieval de precedentes. La evaluación usa 250 consultas por split y un índice reproducible de 2500 casos de train.}
\label{fig:retrieval}
\end{figure}

\begin{table}[H]
\centering
\caption{Resultados de retrieval. Fuente: tabla final en artifacts/metrics.}
\label{tab:retrieval}
\small
\begin{tabular}{@{}llrrrr@{}}
\toprule
\textbf{Método} & \textbf{Split} & \textbf{Consultas} & \textbf{nDCG@10} & \textbf{Recall@5} & \textbf{Recall@10}\\
\midrule
TF-IDF coseno sampled & validation & 250 & 0.883 & 0.948 & 0.964\\
TF-IDF coseno sampled & test & 250 & 0.868 & 0.940 & 0.968\\
TF-IDF SVD coseno sampled & validation & 250 & 0.892 & 0.956 & 0.960\\
TF-IDF SVD coseno sampled & test & 250 & 0.870 & 0.928 & 0.956\\
\bottomrule
\end{tabular}
\end{table}

\subsection{Explicabilidad}

La Tabla~\ref{tab:xai} resume los artefactos XAI. Se extraen términos globales para 10 etiquetas y se generan 5 explicaciones locales. La Tabla~\ref{tab:xai_terms} resume los términos positivos y negativos más influyentes por artículo, y la Figura~\ref{fig:xai} muestra un ejemplo local LIME junto con una prueba mínima de fidelidad. El objetivo es auditar señales del modelo, no demostrar razonamiento jurídico.

\begin{table}[H]
\centering
\caption{Resumen XAI. Fuente: tabla final en artifacts/metrics.}
\label{tab:xai}
\small
\begin{tabular}{@{}rrrr@{}}
\toprule
\textbf{Etiquetas globales} & \textbf{Explicaciones locales} & \textbf{Peso global abs. medio} & \textbf{Prob. local positiva media}\\
\midrule
10 & 5 & 1.372 & 0.736\\
\bottomrule
\end{tabular}
\end{table}

\begin{table}[H]
\centering
\caption{Términos globales principales por artículo. Fuente: coeficientes del modelo lineal final.}
\label{tab:xai_terms}
\scriptsize
\resizebox{\textwidth}{!}{%
\begin{tabular}{@{}lll@{}}
\toprule
\textbf{Artículo} & \textbf{Términos positivos} & \textbf{Términos negativos}\\
\midrule
2 & death (3.36); the death (2.10); son (1.96) & the applicant (-1.52); applicant (-1.15); applicant had (-1.11)\\
3 & cell (3.42); treatment (2.86); medical (2.50) & court of (-1.15); cassation (-1.08); hearing (-1.06)\\
5 & detention (6.72); release (4.61); applicant detention (4.10) & judgment (-1.09); claim (-0.99); 1998 (-0.95)\\
6 & hearing (3.40); hearings (2.25); the proceedings (2.22) & detention (-3.58); applicant detention (-1.92); release (-1.89)\\
8 & child (3.07); family (2.61); censored (2.45) & hearing (-1.19); the judgment (-1.01); hearings (-0.93)\\
9 & religious (3.83); the religious (1.52); of religious (1.48) & regional court (-0.47); compensation (-0.45); 2004 (-0.41)\\
10 & newspaper (3.01); published (2.67); the article (2.46) & detention (-0.90); 2006 (-0.72); the applicant (-0.72)\\
11 & demonstration (2.93); association (1.86); the demonstration (1.84) & the applicant (-0.86); applicant (-0.71); proceedings (-0.71)\\
14 & discrimination (1.79); roma (1.62); religious (1.45) & district court (-1.00); 2004 (-0.98); judgment (-0.90)\\
P1-1 & land (2.61); enforcement (2.27); property (2.24) & hearing (-1.50); hearings (-1.36); detention (-1.24)\\
\bottomrule
\end{tabular}
}
\end{table}

\begin{figure}[H]
\centering
\includegraphics[width=0.88\textwidth]{fig06_xai_global_local.pdf}
\caption{XAI local. LIME y la prueba de retirada de término se interpretan como señales auditables del modelo, no como causalidad jurídica ni justificación doctrinal.}
\label{fig:xai}
\end{figure}

\subsection{Drift y patrones de error}

La Tabla~\ref{tab:drift} muestra un shift moderado de etiquetas. En test, la divergencia JS respecto a train es 0.026 y la similitud coseno 0.951. Macro-F1 aumenta ligeramente respecto a validation, mientras micro-F1 baja y Hamming loss aumenta. Esto sugiere que el comportamiento no se resume con una sola métrica.

\begin{table}[H]
\centering
\caption{Drift de etiquetas y cambios de rendimiento respecto a validation. Fuente: tabla final en artifacts/metrics.}
\label{tab:drift}
\small
\begin{tabular}{@{}lrrrrrr@{}}
\toprule
\textbf{Split} & \textbf{JS} & \textbf{L1} & \textbf{Coseno} & \(\Delta\)\textbf{Macro-F1} & \(\Delta\)\textbf{Micro-F1} & \(\Delta\)\textbf{Hamming}\\
\midrule
validation & 0.017 & 0.278 & 0.959 & 0.000 & 0.000 & 0.000\\
test & 0.026 & 0.320 & 0.951 & +0.018 & -0.022 & +0.007\\
\bottomrule
\end{tabular}
\end{table}

\begin{figure}[H]
\centering
\includegraphics[width=\textwidth]{fig07_drift_error_distributions.pdf}
\caption{Distribuciones de error y label shift. Las líneas punteadas indican medias por split.}
\label{fig:drift_errors}
\end{figure}

La Tabla~\ref{tab:error_patterns} recoge los patrones más frecuentes en test. El caso ideal, sin falsos positivos ni falsos negativos, aparece en 463 casos. Aun así, 194 casos tienen un falso positivo y 165 un falso negativo. Esta diferencia es práctica: un falso positivo suele añadir trabajo de revisión, mientras que un falso negativo puede ocultar una línea argumental relevante.

\begin{table}[H]
\centering
\caption{Patrones de error más frecuentes en test. Fuente: tabla final en artifacts/metrics.}
\label{tab:error_patterns}
\small
\begin{tabular}{@{}rrr@{}}
\toprule
\textbf{FP} & \textbf{FN} & \textbf{Casos}\\
\midrule
0 & 0 & 463\\
1 & 0 & 194\\
0 & 1 & 165\\
1 & 1 & 82\\
0 & 2 & 35\\
2 & 0 & 30\\
\bottomrule
\end{tabular}
\end{table}

\begin{table}[H]
\centering
\caption{Promedios de error por split. Fuente: tabla final en artifacts/metrics.}
\label{tab:error_split}
\small
\begin{tabular}{@{}lrrrr@{}}
\toprule
\textbf{Split} & \textbf{TP} & \textbf{FP} & \textbf{FN} & \textbf{Errores}\\
\midrule
train & 1.325 & 0.237 & 0.138 & 0.374\\
validation & 1.075 & 0.369 & 0.316 & 0.685\\
test & 1.058 & 0.379 & 0.377 & 0.756\\
\bottomrule
\end{tabular}
\end{table}

\begin{figure}[H]
\centering
\includegraphics[width=\textwidth]{fig08_contributions_summary.pdf}
\caption{Mapa compacto de contribuciones: dataset completo, modelo multietiqueta, retrieval, XAI, drift y análisis crítico.}
\label{fig:contrib}
\end{figure}

\section{Discusión}
\label{sec:discusion}

El resultado principal es que un pipeline clásico, reproducible y auditable ofrece una base razonable para priorizar artículos y precedentes en ECtHR Task B. Sin embargo, las métricas no deben confundirse con competencia jurídica. Un macro-F1 de 0.672 en test indica utilidad parcial, no fiabilidad suficiente para operar sin supervisión. En derecho, el coste de un error depende de su tipo y contexto.

Los falsos negativos son especialmente delicados: si el sistema no sugiere un artículo relevante, el usuario puede omitir una línea argumental. Esto puede afectar la estrategia de análisis incluso cuando el resto de etiquetas sean correctas. Los falsos positivos tienen otro tipo de coste: aumentan carga de revisión, pueden introducir ruido y pueden sesgar la búsqueda hacia artículos menos pertinentes. Por ello, en un despliegue real los umbrales deberían ajustarse según coste de error y perfil de usuario, no solo por F1.

El módulo de retrieval puede ahorrar tiempo al localizar casos similares, pero también puede anclar al usuario en precedentes históricos. Si el ranking recupera casos por similitud superficial o por artículos compartidos, el analista podría sobrevalorar analogías débiles. La evaluación utilizada es reproducible, pero limitada: compartir un artículo verdadero no garantiza similitud doctrinal ni relevancia argumentativa.

Las explicaciones por coeficientes y LIME son útiles para auditoría, pero no equivalen a razonamiento jurídico. Un término como ``detention'', ``hearing'' o ``newspaper'' puede ser predictivo de ciertos artículos, pero la relevancia jurídica exige contexto, hechos, estándares doctrinales y valoración normativa. Las explicaciones deben presentarse como señales de inspección, no como justificaciones suficientes.

El drift observado es moderado, pero conceptualmente importante. En un dominio jurídico, un cambio de distribución puede deberse a evolución jurisprudencial, cambios en tipos de demanda o diferencias temporales del benchmark. El sistema debería monitorizar drift, rendimiento y errores de forma periódica, y cualquier política de reentrenamiento debería revisarse con expertos jurídicos.

Finalmente, la privacidad no es el eje experimental del proyecto porque el dataset es público. Aun así, un sistema real aplicado a expedientes o consultas jurídicas podría procesar información sensible. Un despliegue responsable requeriría minimización de datos, anonimización cuando sea posible, control de acceso, registros de auditoría y claridad sobre el uso de modelos externos.

\section{Conclusiones}
\label{sec:conclusiones}

Este trabajo presenta un sistema de apoyo al análisis inicial de casos ante el TEDH basado en recuperación de precedentes, clasificación multietiqueta explicable y análisis de drift. La contribución técnica es un pipeline reproducible que parte de SQLite, entrena TF-IDF + One-vs-Rest, ajusta umbrales, genera predicciones auditables y produce tablas y figuras listas para informe.

Los resultados muestran un comportamiento útil pero limitado. En test, el modelo alcanza macro-F1 de 0.672 y micro-F1 de 0.737, mientras que el retrieval obtiene Recall@10 alto en una evaluación muestral. Estas cifras sugieren que el sistema puede ayudar a priorizar revisión, pero no sostienen automatización de decisiones ni reemplazo de interpretación jurídica.

La lección principal es metodológica: en ML aplicado al derecho, el valor no está solo en mejorar una métrica, sino en conectar resultados con consecuencias reales. Falsos negativos, falsos positivos, anclaje por precedentes, límites de XAI y drift temporal deben discutirse explícitamente. Esta discusión es parte del sistema, no un añadido posterior.

Como trabajo futuro, se propone evaluación cualitativa de precedentes por expertos, análisis temporal por ventanas más finas, calibración de probabilidades, políticas de umbral por coste de error, comparación controlada con modelos jurídicos contextualizados y un diseño de despliegue que incorpore privacidad, auditoría y supervisión humana obligatoria.

\appendix
\section{Anexos}

\subsection{Artefactos reproducibles}

Las tablas finales se guardan en \path{artifacts/metrics/}: \path{paper_classification_table.csv}, \path{paper_retrieval_table.csv}, \path{paper_xai_table.csv}, \path{paper_drift_table.csv}, \path{paper_error_pattern_table.csv} y \path{paper_error_split_table.csv}. Las figuras científicas se generan desde celdas incluidas en los notebooks y se guardan en \path{artifacts/figures/} como PDF y PNG.

\subsection{Uso de Image Gen}

Las Figuras~\ref{fig:overview}--\ref{fig:contrib} son figuras científicas generadas con Matplotlib a partir de datos, métricas o esquemas definidos en los notebooks. De forma separada, se prepararon ilustraciones conceptuales suplementarias con Image Gen para material visual de apoyo. Estas imágenes no sustituyen resultados experimentales ni se usan para inferir métricas. Los prompts se guardan en \path{artifacts/figures/generated/imagegen_prompts.md}. Los cinco prompts documentados fueron:

\begin{enumerate}
\item Ilustración académica limpia de un texto jurídico que fluye hacia vectores TF-IDF, predicción multietiqueta, retrieval de precedentes y explicabilidad, con etiquetas en español.
\item Diagrama de transición de estados: caso ECtHR bruto, estructura SQLite/TF-IDF/scores y salida con artículos candidatos, precedentes y explicación.
\item Figura editorial sobre falsos positivos y falsos negativos en identificación de cuestiones jurídicas.
\item Ilustración conceptual de drift temporal en ML jurídico, con cambio de distribución y política de reentrenamiento.
\item Visual abstract del proyecto: ECtHR Task B, 11000 casos, TF-IDF + One-vs-Rest, retrieval, XAI y drift.
\end{enumerate}

\subsection{Notebooks}

El pipeline final consta de cinco notebooks: \path{00_diagnostico_y_decisiones.ipynb}, \path{01_datos_y_eda.ipynb}, \path{02_modelado_multilabel.ipynb}, \path{03_retrieval_de_precedentes.ipynb} y \path{04_xai_drift_y_errores.ipynb}. El notebook \path{05_guion_informe_y_defensa.ipynb} no forma parte del pipeline final; su contenido útil fue trasladado al README y a este informe.

\begin{thebibliography}{9}

\bibitem{chalkidis2022lexglue}
Chalkidis, I., Jana, A., Hartung, D., Bommarito, M., Androutsopoulos, I., Katz, D. M. y Aletras, N.
\newblock LexGLUE: A benchmark dataset for legal language understanding in English.
\newblock \textit{Proceedings of the 60th Annual Meeting of the Association for Computational Linguistics}, 2022.

\bibitem{chalkidis2020legalbert}
Chalkidis, I., Fergadiotis, M., Malakasiotis, P., Aletras, N. y Androutsopoulos, I.
\newblock LEGAL-BERT: The muppets straight out of law school.
\newblock \textit{Findings of EMNLP}, 2020.

\bibitem{aletras2016predicting}
Aletras, N., Tsarapatsanis, D., Preoţiuc-Pietro, D. y Lampos, V.
\newblock Predicting judicial decisions of the European Court of Human Rights: a natural language processing perspective.
\newblock \textit{PeerJ Computer Science}, 2:e93, 2016.

\bibitem{ribeiro2016lime}
Ribeiro, M. T., Singh, S. y Guestrin, C.
\newblock ``Why should I trust you?'': Explaining the predictions of any classifier.
\newblock \textit{Proceedings of KDD}, 2016.

\bibitem{gama2014survey}
Gama, J., Žliobaitė, I., Bifet, A., Pechenizkiy, M. y Bouchachia, A.
\newblock A survey on concept drift adaptation.
\newblock \textit{ACM Computing Surveys}, 46(4), 2014.

\bibitem{moreno2012datasetshift}
Moreno-Torres, J. G., Raeder, T., Alaiz-Rodríguez, R., Chawla, N. V. y Herrera, F.
\newblock A unifying view on dataset shift in classification.
\newblock \textit{Pattern Recognition}, 45(1):521--530, 2012.

\bibitem{shojaee2025illusion}
Shojaee, P. et al.
\newblock The illusion of thinking: Understanding the strengths and limitations of reasoning models via the lens of problem complexity.
\newblock Apple Machine Learning Research, 2025.

\end{thebibliography}

\end{document}
